\scanpage{19}

On a alors le diagramme commutatif

\diagram{\begin{tikzcd}[ampersand replacement=\&]
\Gamma(V,F) \arrow[r,"\rho'"] \arrow[dr,swap,"\rho"] \& \Gamma(V\cap U',F) \arrow[d,"\rho''"] \\
\& \Gamma(V\cap U,F)=1
\end{tikzcd}}

Si $F$ satisfait à $P_1$ ($\deleted{\unsure{\mathrm{rel}}}P_2$) rel. à $Y$, alors $\rho$ et $\rho''$ sont \deleted{\unsure{des}} inj (bijectifs) puisque $V\cap U=(V\cap U')\cap U$, donc \deleted{leur composé $\rho$ l'est, ce qui} $\rho'$ est inj. (bijectif) ce qui prouve le corollaire.

\deleted{\uline{\unsure{Dém. de prop.} 3 pour $P_1$}}

\uline{Lemme} 5\quad On a, pour tt $x\in X$, \deleted{\unsure{désignant}} \inserted{et tt ouvert $V$,} \deleted{\unsure{posant} $\unsure{U}\ \mathrm{Loc}'_x(X)=\mathrm{Loc}_x(X)-\{x\}$,} $X_{(x)}=\mathrm{loc}_{\deleted{\unsure{x}}x}(X)$, \deleted{$X'_x=X_x-\{x\}$} $\deleted{\unsure{U}}V_{(x)}=\deleted{\unsure{U}}V\cap X_{(x)}$, l'isom. \unsure{suivant}

\[
\deleted{\Gamma(X_x,F)\simeq F_x,\quad \Gamma(\unsure{X'_x|\check X'_x}}
\]

\[
\Gamma(V_{(x)},F_{(x)})\xleftarrow{\sim}
\underset{\substack{W\ \text{voisinage}\\\text{ouvert de }x}}{\varinjlim}\ \Gamma(W\cap V,F)
\tag*{$(*)$}
\]

On a un hom évident, prouvons qu'il est bijectif. Soit $f,g\in\Gamma(W\cap V,F)$ qui \unsure{devient|deviennent} \deleted{\unsure{nulle|nuls}} \inserted{égaux} sur $V_{(x)}$, donc ils sont \unsure{nulle|nuls|nulles} dans un voisinage ouvert \inserted{$E$} de $V_{(x)}\subset X$, \deleted{soit $\unsure{V'_x|W'_x}\ \unsure{E,|E'}$} je dis que $E$ contient un ens. de la forme $W'\cap V$, $W'$ vois. ouvert de $x$.

% archive sheet number: 18
